\chapter{Gebruikersinterface en visualisatie}
\label{chap:interface}

\section{Ontwerp vanuit een operationele schets}
De definitieve interface is gebaseerd op een door de opdrachtgever aangeleverde schets. De nadruk ligt op informatiedichtheid en snelle herkenning, niet op een marketingachtige presentatie. Bovenaan staat een compacte infobalk met sessie, resterende tijd, wagen, huidige piloot, klassepositie, status en acties. Grote vakken tonen laatste en beste rondetijd. Daaronder staan actuele, referentie- en beste sectoren. Rechts krijgen voorspelling, ideal time en de geselecteerde vergelijkingswagen een vaste plaats.

De onderste vergelijkingsmatrix groepeert waarden per betekenis. Elke kolom bevat twee bronwaarden en een delta. Positieve en negatieve uitkomsten krijgen een groene of rode rand volgens hun praktische betekenis. Delta's worden compact zonder overbodige minuten weergegeven, terwijl grote tijdverschillen in de pitstopprojectie wel als minuten en seconden verschijnen.

\section{Race, practice en qualifying}
De modus wordt handmatig in Setup gekozen. Automatische detectie uit de sessienaam werd bewust niet gebruikt, omdat providers sessies anders benoemen en een fout gekozen modus operationeel verwarrend kan zijn.

Race toont het pitstopvenster, recente tempoverschillen en catchindicaties. Practice verbergt racegebonden pitinformatie en focust op referenties, sectoren en vergelijkingen. Qualifying gebruikt beste en laatste tijden in plaats van langetermijngemiddelden. Dezelfde interface blijft herkenbaar, maar de betekenisvolle data verandert per sessietype.

\section{Meerdere gevolgde wagens}
In Setup staat naast het primaire autonummer een plusknop. Hiermee kunnen maximaal twee extra wagens worden toegevoegd. Elke wagen krijgt een eigen dashboardvenster, maar gebruikt dezelfde live snapshot en opslag. Dit is belangrijk voor een team dat tijdens een 24-uurswedstrijd meerdere wagens inzet: de website wordt niet extra belast en alle schermen zijn onderling consistent.

\section{Grafiekvenster}
Via een compact grafiekicoon in de bovenbalk opent een afzonderlijk analysevenster. Vier grafieken kunnen gelijktijdig worden getoond en elk vak heeft een dropdown om het grafiektype te wijzigen. De standaardselectie omvat:

\begin{itemize}
  \item geldige rondetijden per piloot;
  \item beste, gemiddelde en recente prestaties per piloot;
  \item gemiddelde en beste sectoren per piloot;
  \item werkelijke rondetijden van alle wagens in dezelfde klasse.
\end{itemize}

Bij rondetijden per piloot wordt de geldige-rondeteller gebruikt. De derde geldige ronde van piloot A en de derde geldige ronde van piloot B staan daardoor op dezelfde x-positie, ook wanneer hun stints op verschillende momenten plaatsvonden. FCY-, Safety Car-, in- en outlaps laten geen lege plaats achter. Voor grafieken met veel ronden is een zoombereik voorzien, zodat een deel van de wedstrijd kan worden onderzocht zonder de volledige reeks onleesbaar samen te drukken.

\section{Thema en visuele status}
De applicatie ondersteunt een licht en donker thema via een zon- of maanicoon. Kleuren zijn gegroepeerd in CSS-variabelen, zodat het palet later zonder componentwijzigingen kan worden aangepast. De donkere variant gebruikt geen zuiver zwart; panelen blijven visueel van elkaar onderscheiden en rood, oranje en groen behouden voldoende contrast.

Kleur heeft een functionele betekenis. De predictionkaart gebruikt niet alleen een gekleurde rand maar een lichtere volledige achtergrond, zodat een kritische referentietijd direct opvalt. De pitstopcontainer krijgt groen wanneer stoppen momenteel toegestaan is en rood wanneer het venster gesloten is. Tekst blijft daarnaast altijd aanwezig, zodat de toestand niet uitsluitend van kleurwaarneming afhangt.

\section{Setupvensters}
De algemene setup bevat timing-URL, modus, gevolgde wagens en sessiemap. Een afzonderlijke pitstopsetup bevat vaste race-informatie: totale raceduur, verplichte stops, circuit/pitconfiguratie, reguliere trajectafstand en FCY-snelheid. Deze waarden worden voor de race ingesteld en tijdens actieve collectie vergrendeld. Alleen de verwachte pitduur blijft rechtstreeks op het dashboard wijzigbaar.

Referentietijden worden dichter bij hun gebruik aangepast: elk reference-vak heeft een klein editicoon. Hierdoor hoeft de gebruiker niet naar Setup om tijdens een sessie een regelwaarde bij te stellen.

